\scanpage{9}
% L1-3
\deleted{\uline{Remarque} Comme $\alpha_i+2$ est inv i.e. $\alpha_{is}\ne-2$ en tt $s$, cela implique qu'en fait, si \deleted{car.} $\unsure{\nu_i|v_i}=2p$ ($p$ premier \inserted{impair \unsure{donc}}), \deleted{Alors il n'y a pas de $C$} \deleted{on peut en prédire, alors que}}
% L4-5
\deleted{\deleted{et si $S$ contient un pt $s$ de $C'(\nu_0,\nu_1)$ est \unsure{in} car $p$, \unsure{alors que|les $\psi$ sont inv.} Alors $C'(\nu_0,\nu_1)$ est \uline{vide}}}
% L6-9
\deleted{On voit \uline{directmt} que dans ce cas $C(\nu_0,\nu_1)$ est \uline{également} vide : en effet, un \deleted{op.} élément de $G_s$ (forme de $GP(1)_{k(s)}$) ne peut jamais être d'ordre $2p$. Donc il \deleted{serait}\inserted{est} \unsure{raisonnable|vraisemblable},}
% L10-12
\deleted{pour assurer que $C(\nu_0,\nu_1)\deleted{\unsure{\longrightarrow}}\simeq C'(\nu_0,\nu_1)$ soient \uline{évt.} non vides, d'exiger que $\nu_i=2p$ \inserted{($p$ pr. impair)} implique $p$ inv. sur $S$, donc en fait $2p=\nu_i$ inv --}
% L13-15
\deleted{c'est à dire qu'on exige que $\nu_i$ est inv. sur $S$ \uline{ou} \deleted{de la forme $p$} \uline{premier} (impair car $\nu_i\geq3$).}

% L16
Exemples avec $\Gamma_{\nu_0,\nu_1}$ \uline{fini}

% L17
A) $\nu_0=3,\nu_1=3,\qquad\Gamma_{\nu_0,\nu_1}\simeq\unsure{\mathfrak S_4|\mathcal S_4}$. On \deleted{trouve que}\inserted{n'a pas de}
% L18
\deleted{la catégorie des repr. de $\mathfrak S_4$ \unsure{en question}}
% L19
restriction à faire sur $S$. On trouve que la
% L20
cat. des couples $(G,\varphi)$, $G$ forme de $GP(1)_S$,
% L21
$\varphi:\mathfrak S_4\longrightarrow G(S)$ tels que $\forall s\in S$, $\rho_{0s}$ et $\rho_{1s}$ soient
% L22
d'ordre $3$ exactement \inserted{et ne commutent pas}, est \deleted{équivalente à la cat}\inserted{vide si \deleted{$2$ n'est}}
% L23
$2$ ($=(-1)(-1)(-1-1)$ au signe près) n'est pas inversible,
% L24
\deleted{alors} et qu'elle est équivalente à la cat. ponctuelle
% L25
dans le cas contraire. C'est à dire : Si

% L26
a) \deleted{En}\inserted{Sur un corps de} car. $2$, toute une repr. $\mathfrak S_4\unsure{\mathbin{-}|\to}G(S)$, les
% L27
\inserted{$\rho_0$ et $\rho_1$ commutent i.e.} \deleted{$\rho_0$ et $\rho_1$ ou $\rho_1=1$} \deleted{\unsure{i.e. si $\rho_0$ et $\rho_1$ commutent|$\rho_0$ et $\rho_1$ commutent}} \deleted{i.e. $\tau_0\deleted{\unsure{\pm|=}}\tau_1$ ou $\tau_1=\tau_2$ alors} la
% L28
repr. se factorise par $\deleted{\mathfrak S_3(=\mathbb D_3)}\ \deleted{\unsure{\mathfrak S_3}}\ \mathfrak S_3\deleted{\unsure{=\mathbb D_3\subset\mathfrak S_4}}$
% slanted
% L27a
\margin{\deleted{faux?}\ $\big|$}
% archive sheet number: 8
